% Chapter 5, section 1: Run a cluster with Docker
% Write in plain LaTeX. Code of a guide step: \guide{<step>}; commands: \begin{shell}...\end{shell};
% a screenshot: \screenshot{file}{caption}; a gap in the draft: \placeholder{...}.
% See book/latex/queen.sty for the full list.
%
% Draft by Claude, 2026-10-23. Every command below ran on 2.6.0: the Compose cluster on
% other ports (51632..53632), the docker run form on three containers with fixed IPs on one
% Docker network (10.77.0.11..13 instead of 10.0.0.11..13). Outputs are copied from those runs.

\section{Run a cluster with Docker}

A cluster is three copies of the same broker, each on its own machine, each told where the other
two are. Together they keep one log. The leader puts every write in that log and sends it to the
others, and a write is answered only when a majority of the nodes (two of three) has it on disk.
Any node serves any client: it prepares the work of its own clients and sends the leader only the
prepared command.

\why{a machine can die, and with it the disk. With three nodes, losing one costs nothing: the
other two keep working, and no message they acknowledged is lost. The price is three machines,
and a write that waits for a second disk.}

\queenfig{cluster-write}{A write is answered when two of three nodes have it on disk.}

\subsection{Try it on one machine}

The repository has a Compose file that runs three nodes on your laptop, each with its own volume.
Give the nodes a shared secret, start them, and ask each one what it is:

\begin{shell}
cd deploy/compose/three-node        # in a clone of github.com/queen-mq/queen
echo "QUEEN_RAFT_TOKEN=$(openssl rand -hex 32)" > .env
docker compose up -d --wait
for n in 1 2 3; do
  echo "queen-$n $(curl -s localhost:${n}6632/health | jq -r .raft.role)"
done
\end{shell}

\begin{shell}
queen-1 follower
queen-2 follower
queen-3 leader
\end{shell}

The three nodes started empty, voted, and node 3 won. \texttt{-{}-wait} returns when every node is
healthy, a few seconds on a laptop. Node $N$ answers on port $N$6632 (16632, 26632, 36632),
API and dashboard alike. The raft port, 7400, where the nodes talk to each other, is not published
at all.

\queenfig{compose-cluster}{The Compose cluster on one machine.}

Now give a client all three addresses. It spreads its requests over the nodes, and when one does
not answer it moves to the next:

\begin{code}
const queen = new Queen({
  urls: ['http://localhost:16632', 'http://localhost:26632', 'http://localhost:36632'],
})
\end{code}

With node 1 stopped, ten pushes through this client took 64~ms in total, and a consumer read all
twenty messages, the ten pushed before the stop included.

\subsection{Stop the leader}

\begin{shell}
docker compose stop queen-3          # the node that said leader
for n in 1 2; do
  echo "queen-$n $(curl -s localhost:${n}6632/health | jq -r .raft.role)"
done
\end{shell}

\begin{shell}
queen-1 follower
queen-2 leader
\end{shell}

Node 2 leads now. Node 3's log says how long the change took:

\begin{shell}
INFO rsm: raft: leadership handed off to=2 ms=28
\end{shell}

Start node 3 again. It gets the writes it missed from the new leader:

\begin{shell}
docker compose up -d --wait queen-3
curl -s localhost:36632/health | jq -c '.raft | {role, applied, commit, lag}'
\end{shell}

\begin{shell}
{"role":"follower","applied":56,"commit":56,"lag":0}
\end{shell}

\texttt{applied} equal to \texttt{commit} means node 3 has caught up. \texttt{docker compose down -v}
removes the containers and their volumes when you are done.

\why{a planned stop should cost nothing. When a leader gets SIGTERM, it hands its job to the most
up-to-date follower before it exits, so clients see no gap. A crash (kill -9, a lost machine)
costs more: the others wait one to two seconds, then elect a new leader, at most about four
seconds in our tests.}

\queenfig{failover}{The leader stops: a planned hand-off, or an election after a crash.}

\subsection{Three machines}

Three containers on one laptop all die with the laptop, so the Compose file is only for trying
things. In production each node gets its own machine, on a private network. On each machine, run
the same command with that machine's address and node number. Here are the commands for node 1:

\begin{shell}
TOKEN=...              # openssl rand -hex 32, made once: the same secret on every node
PEERS="1=10.0.0.11:7400/10.0.0.11:6632"
PEERS="$PEERS,2=10.0.0.12:7400/10.0.0.12:6632"
PEERS="$PEERS,3=10.0.0.13:7400/10.0.0.13:6632"

docker run -d --name queen --hostname queen-1 \
  --restart unless-stopped --stop-timeout 60 \
  -p 10.0.0.11:6632:6632 -p 10.0.0.11:7400:7400 \
  -v /var/lib/queen:/var/lib/queen/raft \
  -e QUEEN_RAFT_REPLICATOR=openraft \
  -e QUEEN_RAFT_NODE_ID=1 \
  -e QUEEN_RAFT_PEERS="$PEERS" \
  -e QUEEN_RAFT_TOKEN="$TOKEN" \
  ghcr.io/queen-mq/queen:2.7.0
\end{shell}

\begin{options}{Setting}{Default}{What it does}
QUEEN\_RAFT\_REPLICATOR & local & \texttt{openraft} makes the node a member of a cluster. \\
QUEEN\_RAFT\_NODE\_ID & 1 & This node's number: 1, 2, 3. A different number on each node. \\
QUEEN\_RAFT\_PEERS & unset & Every node as \texttt{id=raft address/HTTP address}. The same list on every node. \\
QUEEN\_RAFT\_TOKEN & unset & The secret the nodes show each other. Unset, anyone who reaches port 7400 can vote. \\
QUEEN\_RAFT\_DIR & /var/lib/\allowbreak queen/raft & This node's data. One directory per node, never shared. \\
\end{options}

The Docker flags matter too. \texttt{-{}-stop-timeout 60} gives a stopping node time to hand over
its work (Docker's default of 10 seconds can kill it halfway). \texttt{-{}-restart unless-stopped}
brings a node back after a crash, and after the one planned exit a node makes: when it was away so
long that the leader sent it a snapshot, it exits with code 75 and loads the snapshot when it starts
again. \texttt{-{}-hostname} gives the node a stable name in the logs and the dashboard.

\why{the peer list is how a new cluster is born: every node reads the same list on an empty
directory, and they elect a leader. After that first start the list lives in the log, and the
nodes change it only through the membership API (see Recovery).}

\subsection{Five things to know}

\begin{itemize}
\item \textbf{The addresses are fixed at the first start.} They are stored in the log. Changing a
host name or an IP under a running cluster cuts the nodes off from each other: put the old
addresses back, or run a forced recovery. Use addresses that do not change.
\item \textbf{Keep both ports private.} Port 6632 has no authentication unless you turn it on (see
Configuration for different use cases), and port 7400 is plain HTTP with a shared token. Neither
should face the internet.
\item \textbf{One node does not grow into a cluster.} To go from one node to three, start a new
cluster beside the old one and move the clients to it.
\item \textbf{Give every node the same disk size.} A node refuses its own clients' writes at 85\%
full, but writes accepted by the other nodes keep arriving until its disk is full.
\item \textbf{Run NTP on every machine.} Leases are timed on the leader's clock, and they stay
exclusive across a leader change only while the clocks agree within half a second.
\end{itemize}

On the site: \href{https://queenmq.com/manage/docker/}{queenmq.com/manage/docker}.
